% Dokument używa klasy extarticle, ponieważ standardowy article nie obsługuje czcionki 14 pt.
\documentclass[14pt,a5paper,oneside]{extarticle}

% Kodowanie T1 poprawia dzielenie i kopiowanie polskich słów z pliku PDF.
\usepackage[T1]{fontenc}
% UTF-8 pozwala pisać polskie znaki bez specjalnych komend.
\usepackage[utf8]{inputenc}
% Język polski ustawia poprawne nazwy spisu treści i reguły dzielenia wyrazów.
\usepackage[polish]{babel}
% Latin Modern jest czytelną czcionką ekranową zgodną z kodowaniem T1.
\usepackage{lmodern}
% Geometry ustawia bezpieczne marginesy A5, aby czytnik nie ucinał prawej strony.
\usepackage[
  a5paper,
  inner=13mm,
  outer=13mm,
  top=14mm,
  bottom=16mm,
  headheight=16pt,
  includeheadfoot
]{geometry}
% Microtype delikatnie poprawia skład i zmniejsza ryzyko wychodzenia tekstu poza margines.
\usepackage{microtype}
% Setspace pozwala ustawić wygodną interlinię dla ekranu czytnika.
\usepackage{setspace}
% Enumitem zmniejsza zbyt duże odstępy list w małym formacie A5.
\usepackage{enumitem}
% Xcolor udostępnia spokojny kolor linków.
\usepackage{xcolor}
% Xurl pozwala łamać długie adresy internetowe w dowolnym bezpiecznym miejscu.
\usepackage{xurl}
% Hyperref tworzy klikalny spis treści i zakładki widoczne w czytnikach Kindle.
\usepackage[
  unicode=true,
  bookmarks=true,
  bookmarksopen=true,
  bookmarksnumbered=true,
  colorlinks=true,
  linkcolor=black,
  urlcolor=blue!55!black,
  pdfborder={0 0 0}
]{hyperref}
% Fancyhdr dodaje krótki nagłówek oraz numer strony.
\usepackage{fancyhdr}

% Metadane pomagają bibliotece czytnika rozpoznać dokument.
\hypersetup{
  pdftitle={NieZmoknij - opis projektu},
  pdfauthor={Projekt studencki PAI},
  pdfsubject={Budowa edukacyjnej aplikacji internetowej},
  pdfkeywords={Django, React, PostgreSQL, Redis, Celery, Leaflet, API}
}

% Większa tolerancja składu chroni tekst przed wyjściem poza prawy margines.
\setlength{\emergencystretch}{3em}
% Sloppy jest świadomym zabezpieczeniem małej strony A5 przed przepełnionymi wierszami.
\sloppy
% Interlinia 1.08 zachowuje czytelność bez nadmiernego wydłużania dokumentu.
\setstretch{1.08}
% Listy są zwarte, ale nadal czytelne.
\setlist{itemsep=3pt, topsep=5pt, leftmargin=6mm}
% Spis treści obejmuje sekcje i podsekcje.
\setcounter{tocdepth}{2}
% Numery obejmują poziom podsekcji.
\setcounter{secnumdepth}{2}
% Adresy URL używają tej samej czytelnej rodziny pisma.
\urlstyle{same}

% Nagłówek jest krótki, aby mieścił się na stronie A5.
\pagestyle{fancy}
\fancyhf{}
\fancyhead[L]{NieZmoknij}
\fancyfoot[C]{\thepage}
\renewcommand{\headrulewidth}{0.3pt}

% Polecenie wyróżnia angielski termin techniczny, po którym pojawia się polskie wyjaśnienie.
\newcommand{\term}[1]{\textbf{\textit{#1}}}
% Polecenie formatuje krótkie nazwy plików i endpointów bez ryzyka złamania składni.
\newcommand{\code}[1]{\texttt{\detokenize{#1}}}

\begin{document}

% Strona tytułowa nie ma nagłówka, aby pozostała spokojna i czytelna.
\begin{titlepage}
  \thispagestyle{empty}
  \centering
  \vspace*{18mm}
  {\LARGE\bfseries NieZmoknij\par}
  \vspace{5mm}
  {\Large Opis budowy aplikacji internetowej\par}
  \vspace{13mm}
  {\large Materiał edukacyjny dla ucznia\\Projektowania Aplikacji Internetowych\par}
  \vfill
  {\large Wersja: 13 czerwca 2026 r.\par}
\end{titlepage}

% Spis treści jest zwykłą stroną PDF i jednocześnie tworzy zakładki dla czytnika.
\tableofcontents
\clearpage

\section{Co budujemy}

NieZmoknij jest aplikacją internetową zbierającą dane o pogodzie,
trzęsieniach ziemi, burzach, cyklonach i wulkanach. Użytkownik ogląda te dane
na mapie, może zapisać własne lokalizacje i zobaczyć historię pogody.
Administrator może ręcznie uruchamiać pobieranie danych oraz sprawdzać błędy.

Projekt jest edukacyjny. Nie zastępuje systemu ostrzegania ani komunikatów
służb państwowych. Jego głównym celem jest pokazanie, jak współpracują różne
części nowoczesnej aplikacji.

\subsection{Najważniejsze pojęcia}

\term{Frontend} to część widoczna w przeglądarce. Rysuje mapę, formularze,
tabele i wykresy.

\term{Backend} to część serwerowa. Przyjmuje żądania, sprawdza uprawnienia,
pobiera dane z zewnętrznych źródeł i rozmawia z bazą danych.

\term{API}, czyli \textit{Application Programming Interface}, jest umową
między programami. Określa adresy żądań i format odpowiedzi.

\term{Database} to trwała baza danych. Trwała oznacza, że dane nie znikają
po odświeżeniu strony albo ponownym uruchomieniu przeglądarki.

\term{Container} to odizolowane środowisko procesu. Kontener zawiera program
i potrzebne mu zależności, dzięki czemu projekt łatwiej uruchomić na innym
komputerze.

\section{Jak projekt powstawał}

\subsection{Szkielet i podział odpowiedzialności}

Na początku projekt został podzielony na katalog \code{frontend} oraz
\code{backend}. Ten podział jest ważny: interfejs nie powinien samodzielnie
łączyć się z bazą danych, a backend nie powinien rysować mapy w przeglądarce.

Zostały wybrane:

\begin{itemize}
  \item React i Vite dla interfejsu,
  \item Django i Django REST Framework dla API,
  \item PostgreSQL dla trwałych danych,
  \item Redis dla szybkiego cache i kolejki,
  \item Celery dla zadań wykonywanych w tle,
  \item Docker Compose do wspólnego uruchamiania usług.
\end{itemize}

\subsection{Najpierw modele danych}

\term{Model} opisuje strukturę informacji w bazie. Przykładowo
\code{SavedLocation} przechowuje nazwę, kraj i współrzędne. Każda lokalizacja
ma właściciela. \code{WeatherSnapshot} przechowuje pomiar pogody i wskazuje
lokalizację przez \term{foreign key}, czyli klucz obcy.

Najważniejsza relacja wygląda tak:

\begin{center}
\texttt{User -> SavedLocation -> WeatherSnapshot}
\end{center}

Jeden użytkownik może mieć wiele lokalizacji. Jedna lokalizacja może mieć
wiele pomiarów. Taki układ ogranicza powtarzanie tych samych danych.

\term{Migration} to plik opisujący zmianę schematu bazy. Django wykonuje
migracje w kolejności. Dzięki temu inny komputer może odtworzyć dokładnie
taką samą strukturę tabel.

\subsection{API i pierwsze widoki}

Backend otrzymał endpointy, czyli konkretne adresy API. Przykłady:

\begin{itemize}
  \item \code{GET /api/earthquakes/} zwraca trzęsienia ziemi,
  \item \code{GET /api/locations/} zwraca lokalizacje użytkownika,
  \item \code{GET /api/dashboard/summary/} zwraca agregacje,
  \item \code{POST /api/admin/sync/volcanoes/} uruchamia synchronizację.
\end{itemize}

\term{GET} służy głównie do odczytu. \term{POST} tworzy zasób albo uruchamia
operację. \term{DELETE} usuwa zasób. Odpowiedź ma format \term{JSON}, czyli
tekstową strukturę z polami, listami i wartościami.

\section{Technologie i powody wyboru}

\subsection{Django i Django REST Framework}

Django daje \term{ORM}, czyli warstwę pozwalającą pisać zapytania do bazy
za pomocą obiektów Pythona. Udostępnia też migracje, panel administratora,
walidację i system użytkowników.

Django REST Framework dodaje \term{serializer}. Serializer sprawdza dane
wejściowe i zamienia modele na JSON. Dodaje także klasy uprawnień, na przykład
\code{IsAuthenticated} oraz \code{IsAdminUser}.

Zalety:

\begin{itemize}
  \item wiele potrzebnych elementów jest spójnych i gotowych,
  \item łatwo tworzyć testy API,
  \item panel Django pomaga administratorowi oglądać bazę,
  \item OpenAPI można generować z kodu.
\end{itemize}

Alternatywą byłby FastAPI. Jest lekki i ma bardzo dobrą obsługę typów, ale
panel administratora i część struktury należałoby dobrać osobno. Inną
alternatywą jest NestJS w TypeScript. Daje uporządkowaną architekturę, lecz
w tym projekcie Django szybciej zapewnia bazę, migracje i administrację.

\subsection{React i Vite}

React buduje interfejs z \term{component}, czyli małych części takich jak
mapa, panel logowania lub tabela. \term{State} oznacza dane, które zmieniają
się podczas działania widoku. Zmiana stanu powoduje ponowne narysowanie
odpowiedniej części strony.

Vite jest narzędziem uruchamiającym środowisko developerskie i tworzącym
\term{production build}, czyli zoptymalizowany zestaw plików do publikacji.

Zalety:

\begin{itemize}
  \item strona nie musi przeładowywać się po każdym kliknięciu,
  \item łatwo rozdzielić widoki i stany ładowania,
  \item bogaty ekosystem map i wykresów,
  \item szybki start środowiska programistycznego.
\end{itemize}

Alternatywami są Vue i Svelte. Oba byłyby poprawne. Vue ma łagodny próg
wejścia, a Svelte tworzy małe paczki kodu. React został wybrany ze względu
na bibliotekę React Leaflet i doświadczenie potrzebne do dynamicznego
dashboardu.

\subsection{PostgreSQL}

PostgreSQL jest relacyjną bazą danych. \term{Relational database} przechowuje
dane w tabelach i łączy je kluczami. Projekt korzysta z:

\begin{itemize}
  \item \term{constraint}, czyli ograniczeń poprawności,
  \item \term{index}, czyli dodatkowej struktury przyspieszającej wyszukiwanie,
  \item \term{transaction}, czyli grupy zmian wykonywanych wspólnie,
  \item \term{cascade delete}, czyli usuwania zależnych rekordów.
\end{itemize}

Alternatywą jest SQLite. Jest prostszy i świetny do małych testów, ale
PostgreSQL lepiej przedstawia realistyczne środowisko wielousługowe.
MongoDB byłoby elastyczne dla surowych dokumentów API, lecz relacje
użytkownika, lokalizacji i historii naturalnie pasują do SQL.

\subsection{Redis i cache}

\term{Cache} to szybka pamięć przechowująca wynik, który można później
odtworzyć. Redis przechowuje dane w pamięci operacyjnej i jest dostępny dla
backendu oraz workera.

\term{TTL}, czyli \textit{Time To Live}, określa czas życia wpisu. Pogoda
lokalizacji ma TTL 15 minut. W tym czasie kolejne kliknięcie nie wysyła
identycznego żądania do Open-Meteo.

\term{Invalidation} oznacza usunięcie nieaktualnego cache. Po synchronizacji
trzęsień ziemi cache dashboardu jest usuwany, aby następny odczyt policzył
nowe statystyki.

Alternatywą jest pamięć procesu Django. Byłaby prostsza, ale różne procesy
nie widziałyby wspólnego cache, a restart usuwałby wszystkie wpisy.

\subsection{Celery}

Celery jest systemem \term{task queue}, czyli kolejki zadań. Backend dodaje
zadanie do Redisa, a osobny \term{worker} wykonuje je w tle.
\term{Celery Beat} jest schedulerem, czyli procesem uruchamiającym zadania
według harmonogramu.

Korzyść jest prosta: użytkownik nie czeka, aż serwer pobierze tysiące
rekordów. Trzęsienia są synchronizowane często, pogoda co 15 minut, a duży
katalog wulkanów raz dziennie.

Alternatywą byłby systemowy cron. Jest prostszy, lecz gorzej integruje się
z logami zadań i ręcznym uruchamianiem synchronizacji z panelu.

\subsection{Docker Compose}

Docker Compose opisuje kilka kontenerów w jednym pliku. Projekt uruchamia
frontend, backend, PostgreSQL, Redis, worker i scheduler.

\term{Healthcheck} jest krótkim testem sprawdzającym, czy usługa odpowiada.
Backend startuje po bazie i Redisie, a frontend czeka na zdrowy backend.
Zmniejsza to liczbę błędów wynikających wyłącznie ze złej kolejności startu.

\section{Skąd pochodzą dane}

\subsection{Pogoda}

Open-Meteo zwraca temperaturę, wilgotność, ciśnienie, wiatr, kod pogody,
opad i zachmurzenie. Backend buduje adres z parametrami, pobiera JSON,
normalizuje nazwy pól i zwraca spójny format frontendowi.

Globalna mapa obejmuje stolice świata, największe miasta państw G20 i
20 największych miast Polski. Współrzędne stolic pochodzą z World Bank API.
Punkty są łączone i usuwane są duplikaty.

\subsection{Trzęsienia ziemi}

USGS udostępnia GeoJSON. \term{GeoJSON} to odmiana JSON przeznaczona dla
obiektów geograficznych. Współrzędne mają kolejność longitude, latitude,
czyli długość, a potem szerokość.

Zdarzenia są zapisywane przez \code{update_or_create}. Ta operacja aktualizuje
rekord o znanym identyfikatorze zamiast tworzyć duplikat. Frontend filtruje
czas, magnitudę, głębokość i tekst regionu, a właściwe filtrowanie wykonuje
PostgreSQL.

\subsection{Wulkany i VEI}

Smithsonian Global Volcanism Program udostępnia dane przez
\term{WFS}, czyli \textit{Web Feature Service}. Jest to standard pobierania
obiektów geograficznych. Aplikacja pobiera:

\begin{itemize}
  \item katalog 1215 wulkanów holoceńskich,
  \item osobną historię erupcji.
\end{itemize}

Obie listy są łączone przez \code{Volcano_Number}. \term{VEI}, czyli
\textit{Volcanic Explosivity Index}, opisuje siłę konkretnej erupcji w skali
zwykle od 0 do 8. Nie jest stałą cechą wulkanu.

Dlatego aplikacja pokazuje:

\begin{itemize}
  \item VEI ostatniej znanej erupcji,
  \item najwyższe znane VEI w historii wulkanu,
  \item brak danych, gdy źródło nie sklasyfikowało erupcji.
\end{itemize}

Każdy wulkan jest zaznaczony ikoną \code{pictures/Volcano.png}. Panel pokazuje
także kraj, region, typ, wysokość, rok erupcji i opis geologiczny.

Źródło: \url{https://volcano.si.edu/database/webservices.cfm}.

\subsection{Burze i cyklony}

Potencjał burzowy jest obliczany z bieżących danych Open-Meteo, między innymi
opadu, porywów wiatru i kodu pogody. Nie jest to oficjalny alert.

Cyklony pochodzą z NASA EONET. Są oddzielone od punktów potencjału burzowego,
ponieważ cyklon jest zorganizowanym systemem z własną nazwą i historią.

\section{Logowanie i bezpieczeństwo}

Google Identity Services zwraca \term{ID token}. Backend sprawdza jego podpis,
odbiorcę, ważność i potwierdzenie adresu e-mail. Następnie wydaje własny
\term{access token} oraz \term{refresh token} JWT.

\term{JWT}, czyli \textit{JSON Web Token}, jest podpisanym tekstem zawierającym
informację o użytkowniku. Access token jest krótko ważny. Refresh token służy
do uzyskania nowego access tokenu bez ponownego logowania.

Backend zawsze filtruje prywatne lokalizacje po aktualnym użytkowniku.
Samo zmienienie identyfikatora w adresie nie pozwala odczytać cudzego punktu.
Panel synchronizacji dodatkowo wymaga flagi \code{is_staff}.

Wersja edukacyjna przechowuje tokeny w \term{localStorage}. Jest to wygodne,
ale w produkcji warto przenieść refresh token do cookie \code{HttpOnly},
wymusić HTTPS i dodać restrykcyjną \term{Content Security Policy}.

\section{Najważniejsze optymalizacje}

\subsection{Ograniczenie zewnętrznych requestów}

Pogoda jest pobierana paczkami, a nie osobno dla każdego miasta. Redis
przechowuje wynik. Przy awarii wybranych źródeł aplikacja może użyć ostatniej
poprawnej kopii.

\subsection{Idempotentna synchronizacja}

\term{Idempotent} oznacza, że wielokrotne wykonanie operacji daje ten sam
stan końcowy. Stabilne identyfikatory USGS i Smithsonian powodują aktualizację
rekordu zamiast tworzenia kolejnych kopii.

\subsection{Indeksy bazy}

Indeksy dodano do pól często używanych w filtrach: czasu trzęsienia,
magnitudy, regionu, kraju, nazwy wulkanu i maksymalnego VEI. Indeks zajmuje
dodatkowe miejsce, ale przyspiesza odczyt.

\subsection{Lazy loading}

\term{Lazy loading} oznacza ładowanie kodu dopiero wtedy, gdy jest potrzebny.
Widok lokalizacji, dashboard, tabela sejsmiczna i panel synchronizacji są
osobnymi paczkami. Użytkownik oglądający tylko mapę nie musi od razu pobierać
całego kodu wykresów.

\subsection{Ochrona przed niepełnym importem}

GeoServer podaje liczbę wszystkich dopasowanych rekordów. Import porównuje ją
z liczbą odebranych elementów. Jeśli odpowiedź została ucięta, synchronizacja
kończy się błędem i nie usuwa poprzedniego katalogu.

\section{Testy i jakość}

Backend ma testy jednostkowe i integracyjne. Sprawdzają między innymi:

\begin{itemize}
  \item logowanie Google i wydawanie JWT,
  \item ochronę prywatnych endpointów,
  \item walidację współrzędnych,
  \item cache pogody,
  \item filtry trzęsień ziemi,
  \item łączenie wulkanów z historią VEI,
  \item uprawnienia administratora,
  \item agregacje dashboardu.
\end{itemize}

ESLint analizuje kod JavaScript, a Vite wykonuje próbny build produkcyjny.
GitHub Actions automatycznie uruchamia testy backendu, lint i build po
wysłaniu zmian do repozytorium.

\term{CI}, czyli \textit{Continuous Integration}, oznacza automatyczne
sprawdzanie zmian. Nie zastępuje przeglądu kodu, ale szybko wykrywa wiele
regresji.

\section{Jak działa pojedyncze kliknięcie}

Przykład: użytkownik wybiera warstwę wulkanów.

\begin{enumerate}
  \item React wywołuje funkcję pobierającą dane.
  \item Axios wysyła GET do endpointu backendu.
  \item Django wykonuje zapytanie przez ORM.
  \item PostgreSQL zwraca katalog wulkanów.
  \item Serializer zamienia rekordy na JSON.
  \item React zapisuje odpowiedź w state.
  \item React Leaflet tworzy markery z ikoną Volcano.png.
  \item Kliknięcie markera zapisuje wybrany wulkan w state.
  \item Panel szczegółów pokazuje VEI i dane geologiczne.
\end{enumerate}

Duże pobieranie ze Smithsonian nie dzieje się podczas tego kliknięcia.
Wykonał je wcześniej worker Celery. To rozdzielenie jest jedną z
najważniejszych decyzji architektonicznych projektu.

\section{Uruchomienie projektu}

Najpierw należy utworzyć lokalne pliki środowiskowe i ustawić identyfikator
Google OAuth. Następnie cały zestaw usług uruchamia polecenie:

\begin{center}
\code{docker compose up --build}
\end{center}

Frontend działa pod adresem \url{http://localhost:5174}, a backend pod
\url{http://localhost:8001}. Dokumentacja Swagger jest dostępna pod
\url{http://localhost:8001/api/docs/}.

\term{Environment variable} to zmienna konfiguracyjna przekazywana programowi
spoza kodu. Sekrety i ustawienia zależne od komputera nie powinny być wpisane
na stałe do repozytorium.

\section{Ograniczenia i dalszy rozwój}

Projekt jest rozbudowany, ale nadal edukacyjny. Najważniejsze ograniczenia:

\begin{itemize}
  \item około 1200 ikon może obciążyć słabszą przeglądarkę,
  \item brak VEI jest częsty i nie wolno go zastępować zgadywaną wartością,
  \item potencjał burzowy nie jest oficjalnym ostrzeżeniem,
  \item token w localStorage wymaga szczególnej ochrony przed XSS,
  \item frontend nie ma jeszcze pełnego zestawu testów komponentów i E2E.
\end{itemize}

Rozsądnym kolejnym krokiem jest \term{marker clustering}, czyli grupowanie
blisko położonych ikon przy małym powiększeniu mapy. Następnie warto dodać
test end-to-end logowania, zapisania lokalizacji i otwarcia historii pogody.

\section{Podsumowanie}

NieZmoknij pokazuje pełny przepływ danych: od zewnętrznego API,
przez zadanie Celery, walidację, PostgreSQL i Redis, aż do mapy React.
Każda główna technologia ma konkretne zadanie. Najważniejszą lekcją nie jest
liczba bibliotek, lecz rozdzielenie odpowiedzialności i świadome opisanie
trade-offów, czyli kosztów każdej decyzji.

\end{document}
